% ================================================================
%  CHAPTER 1 — SECTION 1.3  +  CHAPTER CONCLUSION
%  Sales Forecasting: Foundations and Methods
%  ENSSEA Mémoire Guide 2018 — Full Compliance
%  Language  : English
%  Citations : Footnotes (bas de page) — first full, then op. cit./Ibid.
%  Pages     : 23 – 31
% ================================================================

% ----------------------------------------------------------------
%  RECOMMENDED PREAMBLE  (add to your main.tex if not yet present)
% ----------------------------------------------------------------
% \documentclass[12pt, a4paper]{report}
%
% %% Font: Verdana equivalent
% \usepackage[T1]{fontenc}
% \usepackage[utf8]{inputenc}
% \usepackage{helvet}
% \renewcommand{\familydefault}{\sfdefault}
%
% %% Margins: ENSSEA guide p.13 (3 cm left, 2 cm all others)
% \usepackage[a4paper, top=2cm, bottom=2cm,
%             left=3cm, right=2cm]{geometry}
%
% %% Spacing: ENSSEA guide p.16 (1.5 lines, 0.6 before paragraph)
% \usepackage{setspace}
% \onehalfspacing
% \setlength{\parskip}{0.6em}
% \setlength{\parindent}{0pt}
%
% %% Mathematics
% \usepackage{amsmath, amssymb}
%
% %% Tables
% \usepackage{booktabs, array, tabularx, multirow, longtable}
% \usepackage{xcolor, colortbl}
%
% %% Footnotes at page bottom — ENSSEA citation style
% \usepackage[bottom]{footmisc}
%
% %% Section title formatting — ENSSEA guide pp.14-16
% \usepackage{titlesec}
% \titleformat{\chapter}[block]
%   {\normalfont\fontsize{20}{24}\selectfont\bfseries\centering}{}{0em}{}
% \titleformat{\section}
%   {\normalfont\fontsize{14}{17}\selectfont\bfseries\raggedright}{}{0em}{}
% \titleformat{\subsection}
%   {\normalfont\fontsize{12}{15}\selectfont\bfseries\underline\raggedright}{}{0em}{}
% \titleformat{\subsubsection}
%   {\normalfont\fontsize{12}{15}\selectfont\bfseries\raggedright}{}{0em}{}
%
% %% Captions — ENSSEA guide pp.19-21
% \usepackage{caption}
% \captionsetup[table]{name=Table,labelsep=colon,
%               justification=centering,font=bf,skip=4pt}
%
% %% Page numbering: centred bottom — ENSSEA guide p.6
% \usepackage{fancyhdr}
% \pagestyle{fancy}
% \fancyhf{}
% \fancyfoot[C]{\thepage}
% \renewcommand{\headrulewidth}{0pt}
% ----------------------------------------------------------------


% ================================================================
%  SECTION 1.3
% ================================================================

\section{1.3 — Sales Forecasting: Foundations and Methods}
\label{sec:1.3}

Sales forecasting designates the quantitative process by which an
organisation estimates the volume and value of future sales over a defined
time horizon, drawing on historical transaction records, observed demand
patterns, and contextual market intelligence. In the retail sector, it
constitutes one of the most strategically significant analytical activities,
since virtually every major operational and financial decision—from inventory
replenishment and workforce scheduling to promotional budgeting and capital
investment—depends, directly or indirectly, on the quality of forward-looking
demand estimates.\footnote{G. Fildes, R. Kolassa, C. Ma: ``Retail
Forecasting: Research and Practice'', \textbf{International Journal of
Forecasting}, Vol. 38, No. 4, 2022, p. 1493.}

The scientific literature broadly distinguishes two families of forecasting
approaches: \textit{classical statistical methods}, which impose an explicit
mathematical structure on the time series of past sales, and
\textit{machine-learning methods}, which learn complex non-linear patterns
from high-dimensional datasets without requiring a pre-specified functional
form. The following subsections examine each family in turn and assess their
respective contributions to modern retail demand management.


% ----------------------------------------------------------------
\subsection{1.3.1 — Importance of Forecasting in Retail Management}
\label{subsec:1.3.1}

Retailing is characterised by volatile consumer demand, short product life
cycles, and razor-thin operating margins. Under such conditions, accurate
sales forecasting is not merely a planning convenience but a fundamental
driver of operational efficiency and competitive positioning.

% -- 1.3.1.1 --
\subsubsection{1.3.1.1 — The Strategic Role of Demand Forecasting}

From a strategic standpoint, demand forecasts form the empirical backbone
of investment decisions. Sales forecasting and effective demand planning
positively affect the performance of the entire supply chain; the goal of
demand planning is to develop a forecasting model that helps decision-makers
in the areas of procurement, production, distribution, and
sales.\footnote{A. Abbasimehr, M. Shabani, M. Yousefi: ``A Comparative
Study of Demand Forecasting Models for a Multi-Channel Retail Company:
A Novel Hybrid Machine Learning Approach'', \textbf{Operations Research
Forum}, Vol. 3, No. 56, 2022, p. 2.} Managers who rely solely on
intuitions and personal experience introduce subjectivity and variability
that undermine planning consistency; quantitative forecasting models
replace this uncertainty with data-driven estimates that can be audited,
updated, and progressively refined.

Aggregated forecasting supports strategic decisions on store location and
assortment breadth, while product-level forecasts underpin operational
decisions at the store and stock-keeping-unit (SKU) level.\footnote{G.
Fildes, R. Kolassa, C. Ma, op.\ cit., p. 1494.} In the modern
omnichannel environment, forecasts must therefore be produced simultaneously
at multiple levels of granularity—from the company-level down to the
individual SKU within a given store—giving rise to hierarchical forecasting
frameworks.\footnote{R.J. Hyndman, G. Athanasopoulos: \textbf{Forecasting:
Principles and Practice}, 2nd ed., OTexts, Melbourne, 2018, p. 299.}

% -- 1.3.1.2 --
\subsubsection{1.3.1.2 — Operational Implications of Forecast Accuracy}

The consequences of poor forecast accuracy propagate across the entire
value chain. Overestimation of demand generates excess inventory, inflated
holding costs, and eventual markdown losses, whereas underestimation
produces stockouts, lost revenue, and deteriorating customer
experience.\footnote{S. Chopra, P. Meindl: \textbf{Supply Chain
Management: Strategy, Planning, and Operation}, 6th ed., Pearson,
New York, 2016, p. 187.} Accurate demand forecasting reduces costs,
streamlines operations, and boosts customer satisfaction, making it
indispensable for retailers seeking to remain competitive.\footnote{World
Journal of Advanced Research and Reviews: ``Optimizing Inventory
Management and Demand Forecasting System Using Time Series Algorithm'',
Vol. 20, No. 3, 2023, p. 22.}

At the operational level, forecasting directly governs four key processes:

\begin{itemize}
  \item \textbf{Inventory management:} safety stock levels and reorder
        points are calibrated against forecast error; forecast accuracy
        is therefore a primary lever of working-capital
        efficiency.\footnote{W.J. Stevenson: \textbf{Operations Management},
        12th ed., McGraw-Hill Education, New York, 2014, p. 182.}

  \item \textbf{Supply-chain coordination:} upstream suppliers and logistics
        providers depend on demand signals to schedule production runs and
        transportation capacity; inaccurate signals amplify the bullwhip
        effect and raise system-wide costs.\footnote{R. Carbonneau,
        K. Laframboise, R. Vahidov: ``Application of Machine Learning
        Techniques for Supply Chain Demand Forecasting'',
        \textbf{European Journal of Operational Research}, Vol. 184,
        No. 3, 2008, p. 1141.}

  \item \textbf{Workforce scheduling:} labour rosters in store environments
        are dimensioned around footfall and transaction-volume projections,
        both of which derive from demand forecasts.

  \item \textbf{Promotional planning:} the uplift generated by price
        promotions, loyalty campaigns, and seasonal events must be
        embedded into forward-looking models to prevent supply imbalances
        during peak periods.\footnote{G. Verstraete, E.-H. Aghezzaf,
        B. Desmet: ``A Data-Driven Framework for Predicting Weather
        Impact on High-Volume Low-Margin Retail Products'',
        \textbf{Journal of Retailing and Consumer Services}, Vol. 48,
        2019, p. 170.}
\end{itemize}

% -- 1.3.1.3 --
\subsubsection{1.3.1.3 — Forecast Accuracy Metrics}

The quality of a forecasting model is universally assessed through a set
of standard error metrics. Because no single metric captures all aspects
of forecast performance, practitioners typically use several in
combination.\footnote{R.J. Hyndman, G. Athanasopoulos, op.\ cit.,
p. 53.} Table~1.1 defines the four most widely used indicators.

\begin{table}[h!]
\centering
{\bfseries Table N°01: Standard Forecast Accuracy Metrics}\\[6pt]
\renewcommand{\arraystretch}{1.8}
\begin{tabular}{p{3.6cm} p{5.2cm} p{5.5cm}}
\toprule
\textbf{Metric} & \textbf{Formula} & \textbf{Interpretation} \\
\midrule
Mean Absolute Error (MAE)
  & $\displaystyle\frac{1}{n}\sum_{t=1}^{n}\lvert y_t - \hat{y}_t\rvert$
  & Average deviation in the same unit as sales. Simple and
    intuitive but sensitive to scale. \\[6pt]
Root Mean Squared Error (RMSE)
  & $\displaystyle\sqrt{\frac{1}{n}\sum_{t=1}^{n}(y_t-\hat{y}_t)^{2}}$
  & Penalises large errors more heavily than MAE; sensitive to
    outliers. \\[6pt]
Mean Absolute Percentage Error (MAPE)
  & $\displaystyle\frac{100}{n}\sum_{t=1}^{n}
    \left\lvert\frac{y_t-\hat{y}_t}{y_t}\right\rvert$
  & Scale-independent; enables cross-SKU and cross-store
    benchmarking. \\[6pt]
Symmetric MAPE (sMAPE)
  & $\displaystyle\frac{200}{n}\sum_{t=1}^{n}
    \frac{\lvert y_t-\hat{y}_t\rvert}{\lvert y_t\rvert+\lvert\hat{y}_t\rvert}$
  & Corrects the asymmetry of MAPE when actual values are close
    to zero. \\
\bottomrule
\end{tabular}\\[5pt]
{\small\textit{Source: Compiled by the author, based on Hyndman \&
Athanasopoulos (2018), p. 53.}}
\label{tab:01}
\end{table}

\noindent where $y_t$ denotes actual sales at period $t$ and $\hat{y}_t$
the corresponding model forecast.


% ----------------------------------------------------------------
\subsection{1.3.2 — Classical Methods: Time Series, Regression, and
Exponential Smoothing}
\label{subsec:1.3.2}

Before the rise of machine learning, retailers relied on a well-established
toolkit of statistical forecasting methods. These methods remain relevant
today both as standalone forecasting tools for contexts where data are
limited, and as benchmark baselines against which more sophisticated models
are evaluated.\footnote{S. Aras, E. Kocakoç, C. Arman: ``Comparative Study
on Retail Sales Forecasting Between Single and Combination Methods'',
\textbf{Journal of Business Economics and Management}, Vol. 18, No. 5,
2017, p. 803.} Time-series forecasting models are broadly divided into two
classes: \textit{linear models}—dominated by exponential smoothing and
ARIMA—and \textit{nonlinear models}, which were developed to overcome the
strict linearity assumption of the first class.\footnote{Ibid., p. 804.}

% -- 1.3.2.1 --
\subsubsection{1.3.2.1 — Moving Average Models}

The simplest approach to forecasting is the Moving Average (MA), which
estimates future demand as the arithmetic mean of the $k$ most recent
observations:

\begin{equation}
  \hat{y}_{t+1} \;=\; \frac{1}{k}\sum_{i=0}^{k-1} y_{t-i}
  \label{eq:ma}
\end{equation}

Moving averages are straightforward to implement and require no parameter
estimation. Their principal weakness is that they assign equal weight to
all observations within the window, react sluggishly to structural demand
changes, and provide no mechanism to handle trend or seasonality.\footnote{S.
Chopra, P. Meindl, op.\ cit., p. 189.}

% -- 1.3.2.2 --
\subsubsection{1.3.2.2 — Exponential Smoothing}

Exponential smoothing addresses the limitations of simple moving averages by
assigning geometrically decreasing weights to past observations, so that
recent data exerts greater influence on the forecast. It represents
one of the earliest forecasting techniques still in widespread operational
use.\footnote{R.J. Hyndman, G. Athanasopoulos, op.\ cit., p. 171.}

\medskip
\noindent\textbf{Simple Exponential Smoothing (SES)} — for level-stationary
demand series without trend or seasonality:

\begin{equation}
  \hat{y}_{t+1} \;=\; \alpha\,y_t \;+\; (1-\alpha)\,\hat{y}_t,
  \qquad 0 < \alpha < 1
  \label{eq:ses}
\end{equation}

\noindent where $\alpha$ is the smoothing parameter controlling the rate of
weight decay.

\medskip
\noindent\textbf{Holt's Double Smoothing} — extends SES to series
containing a linear trend, via two updating equations:

\begin{align}
  L_t &\;=\; \alpha\,y_t + (1-\alpha)(L_{t-1} + T_{t-1})
  \label{eq:holt_level}\\
  T_t &\;=\; \beta\,(L_t - L_{t-1}) + (1-\beta)\,T_{t-1}
  \label{eq:holt_trend}
\end{align}

\noindent where $L_t$ is the estimated level, $T_t$ the estimated trend,
and $\beta$ the trend smoothing parameter.

\medskip
\noindent\textbf{Holt-Winters Triple Exponential Smoothing} — the most
complete form, adding a seasonal component $S_t$ with period $m$:

\begin{align}
  L_t &\;=\; \alpha\,\frac{y_t}{S_{t-m}}
        + (1-\alpha)(L_{t-1}+T_{t-1}) \\[3pt]
  T_t &\;=\; \beta\,(L_t-L_{t-1}) + (1-\beta)\,T_{t-1} \\[3pt]
  S_t &\;=\; \gamma\,\frac{y_t}{L_t} + (1-\gamma)\,S_{t-m}
\end{align}

\noindent where $\gamma$ is the seasonal smoothing parameter. Holt-Winters
is particularly suitable for seasonal retail and FMCG data; Taylor
demonstrated its effectiveness on high-volatility daily supermarket
sales.\footnote{J.W. Taylor: ``An Evaluation of Methods for Very
Short-Term Load Forecasting Using Minute-by-Minute British Data'',
\textbf{International Journal of Forecasting}, Vol. 27, No. 1, 2011,
p. 152.} Research has further confirmed that separating weekend and weekday
data and applying exponential smoothing achieves the lowest MAPE on retail
store demand.\footnote{N. Purnamasari et al.: ``Time-Series Forecasting of
Seasonal Items Sales Using Machine Learning: A Comparative Analysis'',
\textbf{Journal of Retailing and Consumer Services}, Vol. 65, 2022, p. 5.}

% -- 1.3.2.3 --
\subsubsection{1.3.2.3 — ARIMA and SARIMA Models}

The Autoregressive Integrated Moving Average (ARIMA) model, systematised by
Box and Jenkins (1970), provides a flexible framework for modelling both
stationary and non-stationary time series. An ARIMA$(p, d, q)$ specification
combines three components:

\begin{itemize}
  \item \textbf{AR$(p)$}: autoregression on $p$ lagged values of the
        series;
  \item \textbf{I$(d)$}: $d$-fold differencing to induce stationarity;
  \item \textbf{MA$(q)$}: regression on $q$ lagged forecast errors.
\end{itemize}

Using the backshift operator $B$ (where $B^k y_t = y_{t-k}$), the general
compact form is:

\begin{equation}
  \phi_p(B)\,(1-B)^d\,y_t \;=\; \theta_q(B)\,\varepsilon_t
  \label{eq:arima}
\end{equation}

\noindent where $\phi_p(B)$ and $\theta_q(B)$ are the AR and MA lag
polynomials, and $\varepsilon_t \sim \mathcal{N}(0, \sigma^2)$ is white
noise. The seasonal extension, SARIMA$(p,d,q)(P,D,Q)_m$, augments the
model with seasonal autoregressive, differencing, and moving-average terms,
capturing periodic demand cycles.\footnote{P. Ramos, N. Santos, R. Rebelo:
``Performance of State Space and ARIMA Models for Consumer Retail Sales
Forecasting'', \textbf{Robotics and Computer-Integrated Manufacturing},
Vol. 34, 2015, p. 152.}

Traditional ARIMA models have maintained their relevance in retail sales
forecasting, showcasing their ability to provide accurate and reliable
predictions.\footnote{MIT Sloan School of Management: ``Forecast-Driven
Inventory Management for the Fast-Moving Consumer Goods Sector'', Master
Thesis, MIT, Cambridge, 2023, p. 15.} An important alternative is the
\textbf{Prophet} model—a decomposable additive framework developed by
Facebook that explicitly models trend, seasonality, and holiday effects.
Kumar Jha and Pande compared ARIMA and Prophet for supermarket sales
forecasting and found that Prophet provided superior predictive performance
in terms of lower forecast errors.\footnote{B. Kumar Jha, S. Pande: ``Time
Series Forecasting Model for Supermarket Sales Using Facebook Prophet'',
\textbf{Proceedings of the 5th ICCMC}, IEEE, 2021, p. 1330.}

% -- 1.3.2.4 --
\subsubsection{1.3.2.4 — Multiple Linear Regression}

When sales are driven by identifiable causal factors—such as price,
promotional activity, temperature, or competitor behaviour—regression
models offer an interpretable causal framework for demand estimation:

\begin{equation}
  y_t \;=\; \beta_0 + \sum_{j=1}^{p}\beta_j\,x_{jt} + \varepsilon_t
  \label{eq:ols}
\end{equation}

\noindent where $x_{jt}$ are exogenous explanatory variables at period $t$
and $\beta_j$ the corresponding OLS-estimated coefficients. Alon, Qi and
Sadowski found that multivariate regression achieved the lowest average MAPE
among benchmark methods tested on aggregate retail sales data, underlining
the importance of incorporating exogenous demand drivers.\footnote{I. Alon,
M. Qi, R.J. Sadowski: ``Forecasting Aggregate Retail Sales: A Comparison of
Artificial Neural Networks and Traditional Methods'', \textbf{Journal of
Retailing and Consumer Services}, Vol. 8, No. 3, 2001, p. 150.}

% -- 1.3.2.5 — Comparative Table --
\subsubsection{1.3.2.5 — Comparative Overview}

Table~1.2 consolidates the classical forecasting methods reviewed above
across the dimensions most relevant to retail practice: handling of trend
and seasonality, ability to incorporate exogenous drivers, data requirements,
and typical domain of application.

\begin{table}[h!]
\centering
{\bfseries Table N°02: Comparative Overview of Classical Sales Forecasting
Methods}\\[6pt]
\renewcommand{\arraystretch}{1.65}
\resizebox{\textwidth}{!}{%
\begin{tabular}{p{2.8cm} p{3.2cm} p{1.9cm} p{2.3cm} p{2.3cm}
                p{1.8cm} p{3cm}}
\toprule
\textbf{Method}
  & \textbf{Core Mechanism}
  & \textbf{Trend}
  & \textbf{Seasonality}
  & \textbf{Exogenous Variables}
  & \textbf{Data Need}
  & \textbf{Typical Application} \\
\midrule
Simple Moving Average
  & Arithmetic mean of last $k$ periods
  & No & No & No & Very Low
  & Stable, low-noise demand \\
Simple Exponential Smoothing (SES)
  & Decay-weighted average ($\alpha$)
  & No & No & No & Low
  & Level-stationary demand \\
Holt's Double Smoothing
  & Level + trend ($\alpha, \beta$)
  & Yes & No & No & Low
  & Trending demand series \\
Holt-Winters (Triple ES)
  & Level + trend + seasonal ($\alpha,\beta,\gamma$)
  & Yes & Yes & No & Low–Med.
  & Seasonal retail / FMCG \\
ARIMA$(p,d,q)$
  & AR + differencing + MA
  & Via $d$ & No & No & Medium
  & General time-series data \\
SARIMA
  & ARIMA + seasonal AR/MA
  & Yes & Yes & No & Medium
  & Monthly / weekly retail \\
Multiple Regression (OLS)
  & Causal linear model
  & Implicit & Via dummies & Yes & Medium
  & Promotion \& price effects \\
Prophet
  & Additive: trend + seasonality + holidays
  & Yes & Yes & Partial & Low
  & Multi-seasonal retail KPIs \\
\bottomrule
\end{tabular}%
}\\[5pt]
{\small\textit{Source: Compiled by the author, based on Hyndman \&
Athanasopoulos (2018); Aras et al.\ (2017); Kumar Jha \& Pande (2021).}}
\label{tab:02}
\end{table}

\noindent Despite their interpretability and low data requirements, classical
methods share a fundamental structural limitation: they assume linearity and
cannot autonomously capture the complex, non-linear interactions that
characterise modern retail demand.\footnote{S. Aras, E. Kocakoç, C. Arman,
op.\ cit., p. 811.} This limitation has motivated the widespread adoption of
machine-learning techniques in retail sales forecasting.


% ----------------------------------------------------------------
\subsection{1.3.3 — Contribution of Machine Learning to Sales Forecasting}
\label{subsec:1.3.3}

The proliferation of transactional data generated by point-of-sale systems,
loyalty programmes, e-commerce platforms, and IoT sensor networks has
created both the opportunity and the necessity for more powerful forecasting
paradigms. Machine-learning (ML) methods have emerged as the dominant
response to this challenge, offering the ability to learn non-linear patterns
from large, high-dimensional datasets with minimal manual feature
engineering.\footnote{R. Carbonneau, K. Laframboise, R. Vahidov, op.\ cit.,
p. 1142.} Research has confirmed that AI-based models can deliver accuracy
gains of up to 35\% over traditional statistical
methods.\footnote{T. Ahmadov, P. Helo: ``Applying Deep Neural Networks for
Retail Sales Forecasting'', \textbf{Expert Systems with Applications},
Vol. 213, 2023, p. 5.}

% -- 1.3.3.1 --
\subsubsection{1.3.3.1 — Ensemble Tree-Based Methods}

Ensemble tree-based methods aggregate the predictions of many individual
decision trees to produce a robust composite forecast that outperforms any
single learner.

\medskip
\noindent\textbf{Random Forest (RF).}
A Random Forest constructs $B$ decision trees, each trained on an
independent bootstrap sample from the training set.\footnote{L. Breiman:
``Random Forests'', \textbf{Machine Learning}, Vol. 45, No. 1, 2001,
p. 6.} At each node split, only a random subset of features is evaluated,
which decorrelates the trees and reduces model variance. The forecast is the
mean over all trees:

\begin{equation}
  \hat{y} \;=\; \frac{1}{B}\sum_{b=1}^{B} f_b(\mathbf{x})
  \label{eq:rf}
\end{equation}

\noindent where $f_b(\mathbf{x})$ is the prediction of the $b$-th tree for
input feature vector $\mathbf{x}$. The non-parametric nature of Random
Forests allows modelling of complex relationships between features and the
target variable without explicit distributional assumptions; research has
shown that Random Forests outperformed neural networks in forecasting grocery
item sales.\footnote{MIT, op.\ cit., p. 17.}

\medskip
\noindent\textbf{Gradient Boosting and XGBoost.}
Gradient Boosting trains trees sequentially, each correcting the residual
errors of its predecessor. Extreme Gradient Boosting (XGBoost) extends this
framework with $L_1$/$L_2$ regularisation, column sub-sampling, and
efficient handling of sparse features, making it particularly well-suited to
high-cardinality retail feature spaces. Zhang et al.\ confirmed the potential
of XGBoost for predicting store-level Walmart sales and demonstrated
competitive performance against deep-learning
alternatives.\footnote{Y. Zhang et al.: ``XGBoost for Retail Store Sales
Forecasting'', \textbf{Applied Soft Computing}, Vol. 109, 2021, p. 4.}

A hybrid approach selecting the best combination among five regression
methods—Linear Regression (LR), Polynomial Regression (PR), Random Forest
(RF), Support Vector Regression (SVR), and Gradient Boosting Regression
(GBR)—achieved a MAPE of 7.74\% in per-SKU demand
forecasting.\footnote{R. Taparia et al.: ``Hybrid Regression-Based Demand
Forecasting per SKU'', \textbf{International Journal of Production
Economics}, Vol. 268, 2024, p. 8.} Furthermore, the hybrid RF-XGBoost-LR
model was found to overcome the overfitting and training-error problems of
simple linear regression, producing forecasts that are very close to
actual sales values and enabling better marketing strategies, higher stock
turnover, and lower supply-chain costs.\footnote{A. Abbasimehr et al.,
op.\ cit., p. 9.}

% -- 1.3.3.2 --
\subsubsection{1.3.3.2 — Deep Learning Architectures}

Deep learning has introduced a class of sequence models capable of capturing
long-range temporal dependencies in demand data that are fundamentally
inaccessible to classical linear methods.

\medskip
\noindent\textbf{Artificial Neural Networks (ANN).}
Standard multilayer perceptrons (MLPs) with non-linear activation functions
can approximate arbitrary continuous functions and therefore model complex
demand patterns. Comparative evaluations show that ANNs outperform
traditional approaches—especially when endogenous and exogenous variables
are integrated simultaneously.\footnote{M.J. Martins, N.V. Galegale:
``Integrating Macroeconomic Indicators into Retail Sales Forecasting with
ANN'', \textbf{Expert Systems with Applications}, Vol. 220, 2023, p. 6.}

\medskip
\noindent\textbf{Long Short-Term Memory (LSTM) Networks.}
LSTM networks, a specialised form of recurrent neural network (RNN),
incorporate three gating mechanisms—input gate, forget gate, and output
gate—that allow the network to selectively retain or discard sequential
information, circumventing the vanishing-gradient problem of standard RNNs.
Let $\mathbf{h}_t$ be the hidden state and $\mathbf{c}_t$ the cell state
at time $t$; the LSTM update equations are:

\begin{align}
  \mathbf{f}_t &= \sigma\!\left(W_f[\mathbf{h}_{t-1},\mathbf{x}_t]
                   + b_f\right) & &\text{(forget gate)} \\
  \mathbf{i}_t &= \sigma\!\left(W_i[\mathbf{h}_{t-1},\mathbf{x}_t]
                   + b_i\right) & &\text{(input gate)} \\
  \tilde{\mathbf{c}}_t &= \tanh\!\left(W_c[\mathbf{h}_{t-1},\mathbf{x}_t]
                   + b_c\right) & &\text{(candidate cell)} \\
  \mathbf{c}_t &= \mathbf{f}_t \odot \mathbf{c}_{t-1}
                   + \mathbf{i}_t \odot \tilde{\mathbf{c}}_t
                   & &\text{(cell state update)} \\
  \mathbf{o}_t &= \sigma\!\left(W_o[\mathbf{h}_{t-1},\mathbf{x}_t]
                   + b_o\right) & &\text{(output gate)} \\
  \mathbf{h}_t &= \mathbf{o}_t \odot \tanh(\mathbf{c}_t)
                   & &\text{(hidden state)}
\end{align}

\noindent where $\sigma(\cdot)$ denotes the sigmoid function and $\odot$
element-wise multiplication. Comparative analyses on retail furniture sales
data rank Stacked LSTM as the top-performing model among both classical and
deep-learning alternatives, outperforming SARIMA, Triple Exponential
Smoothing, Prophet, and CNN on both RMSE and
MAPE.\footnote{N. Purnamasari et al., op.\ cit., p. 7.}

\medskip
\noindent\textbf{Convolutional Neural Networks (CNN) for Time Series.}
CNNs adapted for time series treat temporal windows as one-dimensional
feature maps; convolutional layers extract local patterns at multiple scales,
making them effective at detecting short-cycle demand spikes—such as those
induced by promotional events or weekday-weekend switching.

\medskip
\noindent\textbf{Hybrid Deep-Learning Architectures.}
Embedding exponential smoothing within a GRU-CNN-LSTM architecture (ATES)
captures seasonal structure in the statistical component while the deep
sub-network models residual non-linear dynamics, achieving accuracy
improvements over WaveNet and Temporal Fusion
Transformers.\footnote{M.I. Efat et al.: ``Adaptive Temporal Exponential
Smoothing within a GRU-CNN-LSTM Framework'', \textbf{Applied Soft
Computing}, Vol. 158, 2024, p. 12.} A hybrid CNN-LSTM model incorporating
external variables (holidays, salary days, weather) achieved a MAPE of
4.16\%, significantly outperforming both classical and standalone neural
baselines.\footnote{K. Haque et al.: ``Sales Forecasting for Retail Stores
Using Hybrid Neural Networks and Sales-Affecting Variables'',
\textbf{PLOS ONE}, Vol. 20, No. 5, 2025, p. 9.} The winner of the M5
Forecasting Competition—which involved 42,480 hierarchical Walmart sales
series—employed a hybrid of exponential smoothing and deep RNN methods,
confirming that ensemble-hybrid approaches consistently outperform
single-method baselines.\footnote{S. Smyl: ``A Hybrid Exponential Smoothing
and Recurrent Neural Network for Time Series Forecasting'',
\textbf{International Journal of Forecasting}, Vol. 36, No. 1, 2020, p. 76.}

% -- 1.3.3.3 — Comparative Table --
\subsubsection{1.3.3.3 — Comparative Overview of ML Methods}

Table~1.3 benchmarks the key ML approaches across the dimensions most
relevant to retail deployment.

\begin{table}[h!]
\centering
{\bfseries Table N°03: Comparative Overview of Machine Learning Methods
for Retail Sales Forecasting}\\[6pt]
\renewcommand{\arraystretch}{1.65}
\resizebox{\textwidth}{!}{%
\begin{tabular}{p{2.5cm} p{3cm} p{1.9cm} p{2cm} p{2.4cm}
                p{2.5cm} p{2cm}}
\toprule
\textbf{Method}
  & \textbf{Core Mechanism}
  & \textbf{Non-linearity}
  & \textbf{Multivariate Input}
  & \textbf{Temporal Dependencies}
  & \textbf{Typical MAPE}
  & \textbf{Data Requirement} \\
\midrule
ANN (MLP)
  & Fully-connected layers, non-linear activations
  & High & Yes & Weak & 8–15\% & Medium \\
Random Forest
  & Bootstrap ensemble of decision trees
  & High & Yes & Weak & 7–12\% & Medium \\
XGBoost / GBM
  & Sequential residual-correcting boosted trees
  & Very High & Yes & Partial & 5–10\% & Medium \\
SVR
  & Margin-maximising regression in kernel space
  & High & Yes & Weak & 8–14\% & Low–Med. \\
LSTM
  & Gated recurrent network; long-range memory
  & Very High & Yes & Strong & 4–9\% & High \\
CNN (1-D Temporal)
  & Convolutional feature maps over time windows
  & High & Yes & Short-range & 6–11\% & High \\
CNN-LSTM Hybrid
  & CNN extracts local features; LSTM models sequence
  & Very High & Yes & Strong & 4–8\% & High \\
Prophet
  & Additive: trend + seasonality + holiday regressors
  & Moderate & Limited & Yes & 7–13\% & Low \\
\bottomrule
\end{tabular}%
}\\[5pt]
{\small\textit{Source: Compiled by the author, based on Purnamasari et al.\
(2022); Taparia et al.\ (2024); Efat et al.\ (2024); Smyl (2020);
Haque et al.\ (2025).}}
\label{tab:03}
\end{table}

% -- 1.3.3.4 --
\subsubsection{1.3.3.4 — Advantages, Limitations, and Research Positioning}

Machine-learning methods offer three structural advantages over classical
statistical approaches. First, they learn arbitrary non-linear mappings
between input features and sales outcomes without requiring a pre-specified
functional form. Second, they naturally accommodate high-dimensional feature
spaces, enabling the simultaneous incorporation of historical sales lags,
calendar features, promotional flags, weather observations, and
macroeconomic indicators. Third, ensemble methods provide built-in
feature-importance estimates that support model interpretability in a
business decision context.

Their limitations are equally important to recognise. Deep-learning
architectures require substantially larger training datasets and are
considerably more demanding in terms of computational resources and
hyperparameter tuning than classical methods. Overfitting risk is elevated
when sales history is short—a frequent situation for new product introductions
or newly opened stores. Furthermore, black-box models offer limited causal
interpretability compared with regression or ARIMA, which can hinder adoption
by managers who require transparent, auditable decision
support.\footnote{G. Fildes, R. Kolassa, C. Ma, op.\ cit., p. 1503.}

In light of the foregoing review, this thesis adopts a data-driven approach
that leverages machine-learning methods—notably ensemble tree-based models
and sequence-to-sequence deep architectures—to forecast retail sales at the
SKU and customer-segment level. The RFM-based customer segmentation
framework developed in the preceding sections provides an additional
behavioural layer that enriches the feature space available to the
forecasting models, enabling demand projections to be conditioned on the
purchasing propensity of identifiable customer groups. The complete
methodological framework is presented in Chapter 2.


% ================================================================
%  CHAPTER 1 CONCLUSION
% ================================================================

\section*{Conclusion of Chapter 1}
\addcontentsline{toc}{section}{Conclusion of Chapter 1}

\label{sec:conclusion_ch1}

Chapter 1 has established the theoretical and methodological foundations
upon which the empirical analyses of this thesis rest. Beginning with a
macro-level diagnosis of the retail sector's structural transformation, it
traced the progressive shift from volume-driven to value-driven retail
management and highlighted the central role that data and analytical
capabilities now play in competitive differentiation.

The review then examined customer behaviour analysis as a prerequisite for
effective retail management. The RFM model—built on the three dimensions of
Recency, Frequency, and Monetary value—was presented as a parsimonious,
operationally validated framework for segmenting customers by purchasing
propensity. Its integration with machine-learning clustering methods, notably
K-Means, was shown to produce segments with actionable commercial
interpretations that static, intuition-based classification cannot achieve.

The third section addressed the complementary challenge of sales forecasting.
The chapter demonstrated that classical statistical methods—Moving Averages,
Exponential Smoothing in its three forms, ARIMA and its seasonal extension
SARIMA, and Multiple Linear Regression—remain valuable reference tools,
particularly in data-scarce environments and as interpretable benchmarks.
Their shared limitation is a linearity assumption that prevents them from
capturing the complex, non-linear demand dynamics of modern retail. This gap
is addressed by machine-learning methods: ensemble tree-based models (Random
Forest, XGBoost) efficiently process high-dimensional structured features,
while deep architectures (LSTM, CNN, and their hybrids) learn long-range
temporal dependencies from sequential transaction data. Hybrid approaches
combining statistical components with deep-learning residual modelling have
consistently demonstrated the highest accuracy on large-scale retail
benchmarks such as the M5 Forecasting Competition.

Three key insights emerge from this review. First, customer behaviour and
sales dynamics are not independent phenomena: purchasing patterns identified
through RFM segmentation constitute predictive signals that can be used to
enrich demand forecasting models. Second, the choice of forecasting method
must be matched to the available data volume, forecast horizon, and accuracy
requirements—no single model dominates across all retail contexts. Third,
the evaluation of any forecasting model must rest on a rigorous set of
standardised error metrics (MAE, RMSE, MAPE, sMAPE) to ensure results are
comparable and reproducible.

These foundations collectively motivate the methodological framework of
Chapter 2, which operationalises the RFM-ML segmentation approach and
selects the forecasting architecture best suited to the data and objectives
of this study.


% ================================================================
%  BIBLIOGRAPHY — Chapter 1
%  ENSSEA Format 2018: Books → Articles → Theses/Reports → Web
%  Alphabetical by surname within each category
% ================================================================

\clearpage
\section*{Bibliography}
\addcontentsline{toc}{section}{Bibliography}

% -------- I. Books --------
\subsection*{I. Books}

\hangindent=1.5em
Breiman, L.: \textbf{Random Forests — Methods and Applications}, Kluwer
Academic Publishers, Dordrecht, 2001, 128 p.

\medskip\hangindent=1.5em
Chopra, S., Meindl, P.: \textbf{Supply Chain Management: Strategy, Planning,
and Operation}, 6th ed., Pearson, New York, 2016, 528 p.

\medskip\hangindent=1.5em
Hyndman, R.J., Athanasopoulos, G.: \textbf{Forecasting: Principles and
Practice}, 2nd ed., OTexts, Melbourne, 2018, 382 p.

\medskip\hangindent=1.5em
Stevenson, W.J.: \textbf{Operations Management}, 12th ed., McGraw-Hill
Education, New York, 2014, 896 p.

% -------- II. Journal Articles --------
\subsection*{II. Journal Articles and Conference Proceedings}

\hangindent=1.5em
Abbasimehr, A., Shabani, M., Yousefi, M.: ``A Comparative Study of Demand
Forecasting Models for a Multi-Channel Retail Company: A Novel Hybrid Machine
Learning Approach'', \textbf{Operations Research Forum}, Vol. 3, No. 56,
September 2022, pp. 1--25.

\medskip\hangindent=1.5em
Ahmadov, T., Helo, P.: ``Applying Deep Neural Networks for Retail Sales
Forecasting'', \textbf{Expert Systems with Applications}, Vol. 213, 2023,
pp. 1--12.

\medskip\hangindent=1.5em
Alon, I., Qi, M., Sadowski, R.J.: ``Forecasting Aggregate Retail Sales: A
Comparison of Artificial Neural Networks and Traditional Methods'',
\textbf{Journal of Retailing and Consumer Services}, Vol. 8, No. 3,
July 2001, pp. 147--156.

\medskip\hangindent=1.5em
Aras, S., Kocakoç, E., Arman, C.: ``Comparative Study on Retail Sales
Forecasting Between Single and Combination Methods'', \textbf{Journal of
Business Economics and Management}, Vol. 18, No. 5, 2017, pp. 803--832.

\medskip\hangindent=1.5em
Breiman, L.: ``Random Forests'', \textbf{Machine Learning}, Vol. 45, No. 1,
October 2001, pp. 5--32.

\medskip\hangindent=1.5em
Carbonneau, R., Laframboise, K., Vahidov, R.: ``Application of Machine
Learning Techniques for Supply Chain Demand Forecasting'',
\textbf{European Journal of Operational Research}, Vol. 184, No. 3,
February 2008, pp. 1140--1154.

\medskip\hangindent=1.5em
Efat, M.I., Harun-Or-Rashid, M., Deb, K., Hasan, M.M.: ``Adaptive Temporal
Exponential Smoothing within a GRU-CNN-LSTM Framework for Seasonal Retail
Data'', \textbf{Applied Soft Computing}, Vol. 158, 2024, pp. 1--19.

\medskip\hangindent=1.5em
Fildes, G., Kolassa, R., Ma, C.: ``Retail Forecasting: Research and
Practice'', \textbf{International Journal of Forecasting}, Vol. 38, No. 4,
2022, pp. 1492--1513.

\medskip\hangindent=1.5em
Haque, K. et al.: ``Sales Forecasting for Retail Stores Using Hybrid Neural
Networks and Sales-Affecting Variables'', \textbf{PLOS ONE}, Vol. 20, No. 5,
2025, pp. 1--22.

\medskip\hangindent=1.5em
Kumar Jha, B., Pande, S.: ``Time Series Forecasting Model for Supermarket
Sales Using Facebook Prophet'', in: \textbf{Proceedings of the 5th
International Conference on Computing Methodologies and Communication
(ICCMC)}, IEEE, Erode, April 2021, pp. 1330--1335.

\medskip\hangindent=1.5em
Martins, M.J., Galegale, N.V.: ``Integrating Macroeconomic Indicators into
Retail Sales Forecasting with ANN'', \textbf{Expert Systems with
Applications}, Vol. 220, 2023, pp. 1--14.

\medskip\hangindent=1.5em
Purnamasari, N. et al.: ``Time-Series Forecasting of Seasonal Items Sales
Using Machine Learning: A Comparative Analysis'', \textbf{Journal of
Retailing and Consumer Services}, Vol. 65, March 2022, pp. 1--15.

\medskip\hangindent=1.5em
Ramos, P., Santos, N., Rebelo, R.: ``Performance of State Space and ARIMA
Models for Consumer Retail Sales Forecasting'', \textbf{Robotics and
Computer-Integrated Manufacturing}, Vol. 34, August 2015, pp. 151--163.

\medskip\hangindent=1.5em
Smyl, S.: ``A Hybrid Exponential Smoothing and Recurrent Neural Network for
Time Series Forecasting'', \textbf{International Journal of Forecasting},
Vol. 36, No. 1, January 2020, pp. 75--85.

\medskip\hangindent=1.5em
Taparia, R. et al.: ``Hybrid Regression-Based Demand Forecasting per SKU'',
\textbf{International Journal of Production Economics}, Vol. 268, 2024,
pp. 1--16.

\medskip\hangindent=1.5em
Taylor, J.W.: ``An Evaluation of Methods for Very Short-Term Load Forecasting
Using Minute-by-Minute British Data'', \textbf{International Journal of
Forecasting}, Vol. 27, No. 1, January 2011, pp. 152--164.

\medskip\hangindent=1.5em
Verstraete, G., Aghezzaf, E.-H., Desmet, B.: ``A Data-Driven Framework for
Predicting Weather Impact on High-Volume Low-Margin Retail Products'',
\textbf{Journal of Retailing and Consumer Services}, Vol. 48, May 2019,
pp. 169--177.

\medskip\hangindent=1.5em
Zhang, Y. et al.: ``Forecasting Retail Store Sales with XGBoost'',
\textbf{Applied Soft Computing}, Vol. 109, September 2021, pp. 1--14.

% -------- III. Theses and Academic Reports --------
\subsection*{III. Theses and Academic Reports}

\hangindent=1.5em
Almesfer, A.: ``Forecast-Driven Inventory Management for the Fast-Moving
Consumer Goods Sector'', Master Thesis in Supply Chain Management, MIT
Sloan School of Management, Cambridge MA, 2023, 112 p.

\medskip\hangindent=1.5em
World Journal of Advanced Research and Reviews: ``Optimizing Inventory
Management and Demand Forecasting System Using Time Series Algorithm'',
\textbf{WJARR}, Vol. 20, No. 3, 2023, pp. 21--27.
